\chapter{Conclusion\label{cha:chapter7}}


\section{Summary\label{sec:summary}}

Prior research for SIF enhancement has shown that sparse and coarse satellite SIF data can be reconstructed, spatially enhanced (increased resolution) using optical, radiation, meteorological and land cover variables \cite{li2019global,zhang2018global,shen2022spatiotemporal}. Reconstruction methods based mainly on OCO-2 and OCO-3 data give spatial and temporal completeness, but any fine scale values produced are mainly inferred from predictor relationships with SIF and do not usually have a requirement (passed as another predictor) to preserve the coarse SIF observation magnitude from the same time. In contrast TROPOMI based downscaling methods \cite{hu2024spatially,chen2025improved,zhang2023generating} can retain more of the original measured SIF signal by giving its magnitude bounds as a predictor for downscaling the coarse signal, ie; the model will know what is the general range of SIF at the particular pixel it is supposed to downscale thereby reducing the range for prediction. OCO-2 and OCO-3 cannot provide such coarse signal for every point on Earth. These TROPOMI based methods retain the original signal by bias correction, redistribution and physical observation constraints. Reconstruction and Downscaling methods that used a convolutional representation of data showed that spatial predictor grids can represent neighbourhood patterns and mixed land cover patterns that would generally be unavailable to tabular models containing mean aggregated predictor values\cite{gensheimer2022convolutional,du2025cnsif}. However the finer SIF output does not guarantee that the estimated subpixel SIF pattern represents true physiological variation because many of the predictors can represent more about the canopy structure rather than short term photosynthetic regulation, ie; canopy structures changes in a much more slower timeframe than photosynthesis process. This limitation is also more difficult to address with OCO-2 than with TROPOMI because the narrow and widely separated OCO-2 tracks leave large areas without a parent SIF value that could give magnitude bounds to the redistribution method \cite{yu2019high}. Also the predictors can vary slightly in terms of data acquisition and therefore the final fine scale output represents more about the average subpixel SIF distribution than an instantaneous distribution.

The previous highest resolution SIF products also remain at approximately 500m resolution or coarser, which is insufficient for  isolating crop pure SIF signals for many small and irregularly shaped agricultural fields especially in a German setting. Model generalizability and fine scale validation is also a common limitation in the papers shown in Table \ref{tab:sif_enhancement_literature}, as many enhanced SIF products are assessed against held out supervisory observations (same scale, same region, same satellite) or indirect proxies such as GPP and yield rather than independent SIF measured at the output fine resolution. Though the latter is more due to the lack of availability of SIF observations higher than 1km$^2$ resolution, though newer satellite observations like FLEX ESA \cite{drusch2016fluorescence, coppo2017fluorescence} will soon observe 300m resolution SIF data which can then help validate these methods more reliably. Crop yield and productivity studies show that SIF is the most important remote sensing data for the purpose, accurate estimations depends on the crop type, observation timing, geographic scale and inclusion of additional predictors like weather, soil moisture and other environmental predictors \cite{peng2020assessing,song2018satellite,qiu2022monitoring}. The existing SIF resolution limitation and its application to German crop yield prediction motivated the present study to test OCO-2 SIF enhancement with Sentinel-2 predictors at 20m resolution and examine whether the resulting crop pure signal benefits winter wheat yield prediction. 

The experimental workflow for this thesis consisted of two connected stages: SIF enhancement followed by regional winter wheat yield prediction. The first stage used quality controlled OCO-2 observations from February to July during 2019-2024 across nine MGRS tiles in Germany, with two further tiles reserved for external evaluation (model generalizability). SIF quality processing involved combining 757nm and 771nm observations, along with their corresponding noise retrieval uncertainty corrections and checks applied producing a final relatively noise reduced target. Sentinel-2 reflectance bands, spectral indices, GLASS spectral and radiation products, VIIRS radiation data and annual crop type rasters were acquired and matched as closely as possible to the acquisition date of each SIF observation. All data were aligned to a common 20m grid through the process of area averaging, bilinear interpolation. The irregular shaped OCO-2 polygons were retained as fractional masks in the various input spatial supports (image input dimension) so that the predictions could be averaged over their actual spatial support. To understand the effect of model architecture and input data representation, various models like SIF footprint mean aggregated Random Forest and XGBoost models, area-to-point MLP, single and multi footprint convolutional models, and a U-Net trained on fixed 4km spatial aggregates were trained.  These experiments tested whether a spatial model could produce context aware (neighbourhood relationships) fine grid maps and whether OCO-2 aggregation could provide a more stable training target. The resulting 20m SIF outputs should be treated as model derived downscaled estimates rather than direct SIF validated observations (20m), at most it can be validated at 1km$^2$ resolution. While the design allowed to inspect the effects of target construction, predictor representation, spatial support and model family within one workflow, it cannot guarantee a controlled comparison of model architecture alone and should therefore be seen as comparisons of complete modelling configurations. 

Training, validation and testing partitioning was done in such a way that complete dates, footprints, connected date-product components remained together depending on the requirements of each experiment. The best model averaged at least five same date and same mode OCO-2 observations within each fixed 4km grid cell and trained a U-Net whose output was aggregated through the corresponding footprint weight map as shown in Figure \ref{fig:unet}. Internal evaluation used held out aggregate cells while external evaluation was done by applying the saved model and validation derived calibration to native OCO-2 footprints in two previously unseen MGRS tiles. Model performance was additionally examined by month, year, measurement mode, plant-hardiness zone, SIF to Sentinel-2 temporal alignment and footprint count to identify conditions associated with unstable predictions. For the second stage, the selected model was applied to winter wheat pixels in five Bavarian MGRS tiles and mean aggregated from March to July for each NUTS 3 region and year. Random Forest yield models: crop pure enhanced SIF, raw footprints containing at least 5\% winter wheat and all accepted raw footprint model configurations were trained on 2019-2022 year splits and tested on 2024 split. Lastly, paired bootstrap resampling of the 48 test regions was then used to quantify uncertainty in the differences between the crop-pure and raw-SIF models. 

The best model Spatial aggregate (4km) achieved the strongest internal result of RMSE 0.0701 and an R$^2$ 0.761. Aggregation reduced both target and residual variability, confirming that averaging several OCO-2 observations can provide more stable supervision. But it also narrowed the SIF range and removed most extreme values and thus a larger geographic scope is needed to acquire more extreme samples for the model to understand the conditions resulting in very high SIF values. When the model was transferred to individual footprints (1km$^2$) in two unseen MGRS tiles, performance decreased to an RMSE of 0.1434 and an $R^2$ of 0.349. It showed that aggregate accuracy overstates model performance at noisier native support. Regarding model architecture, tree based models produced broadly similar aggregate metrics indicating that convolutional architecture alone does not guarantee much greater predictive accuracy than carefully constructed tabular predictors. The convolutional models nevertheless have an edge when it comes to spatial representation because their 20m SIF maps used neighbouring pixel information and preserved more accurate field boundaries and landscape context. For the yield experiment, crop pure enhanced SIF produced the lowest 2024 RMSE at 0.9406 t/ha, compared with 1.0285 t/ha for the 5\% wheat raw footprint model and 1.0874 t/ha when all raw footprints were used. This corresponded to RMSE reductions of approximately 8.5\% and 13.5\%, which is in line with the general literature stating that more isolated SIF information can provide useful crop yield and productivity information \cite{peng2020assessing,song2018satellite,qiu2022monitoring}. The yield experiment should be interpreted primarily as a comparison of SIF isolation strategy rather than as an assessment of overall yield prediction accuracy given the R$^2$ and 95\% bootstrap intervals reported in Table \ref{tab:yield_model_performance}. Compared to the existing highest resolution enhanced SIF products at 500m resolution \cite{gensheimer2022convolutional,du2025cnsif}, the thesis' best model external validation at 1km$^2$ RMSE = 0.1434 and an $R^2$ = 0.349 is comparable to the performance metrics of SIFnet OCO-2 validated result of RMSE = 0.19, R$^2$ = 0.57, CNSIF's tower based (nine ChinaSpec sites) SIF validation RMSE = 0.152, R$^2$ = 0.581. Though it should be noted exact comparison is not reliable as the existing papers tested on different geographic settings (CONUS, China) and at different spatial scales (Nation wide for SIFnet, small local tower based setting for CNSIF). The yield comparison cannot be done directly with existing papers because of the vast differences in crop type, spatial and temporal setting and should be used more of a general expected range for results, useful predictors and model architecture. The results also agree with SIFnet and CNSIF \cite{gensheimer2022convolutional,du2025cnsif} in showing the value of spatial predictor structure for mixed land cover and also demonstrate that fine output resolution must not be confused with fine scale validation. Finally, the present OCO-2 gap filling and resolution enhancement framework can predict between satellite tracks without a parent measurement unlike the observation preserving TROPOMI methods, but this also increases the importance of the quality of the predictor variables thereby making independent fine resolution validation with future high resolution SIF data like FLEX ESA 300m SIF an important requirement for future work \cite{chen2025improved,zhang2023generating}. 

\section{Dissemination\label{sec:dissemination}}

The main outputs for this thesis are the evaluated SIF enhancement workflows: the associated data preparation and modelling workflow and evidence concerning the usefulness and limitations of crop pure SIF for regional yield prediction. The most immediate potential users are remote-sensing, SIF, and machine-learning researchers who require methods for combining sparse fluorescence observations with spatially complete high resolution predictors. Agricultural researchers and crop monitoring agencies could use the resulting crop specific SIF summaries together with temperature, weather and soil information to investigate crop productivity \& yield and plant stress.  In addition, spatially detailed SIF estimates could support research on vegetation phenology, and plant responses to drought or heat \cite{qiu2022monitoring,song2018satellite,peng2020assessing}. The thesis is itself the main dissemination output and documents the datasets, their quality, processing decisions, model configurations, validation design and interpretation limits needed to understand and reproduce this work. The thesis study should be treated as a research prototype and has not yet been peer reviewed or integrated into an operational monitoring service. Integration into a larger environment would mean connecting the study's data acquisition, preprocessing, model inference to a GIS platform that can support real time inference, data generation and also support external tasks like agricultural monitoring. Such an integration would require a more automated workflow than the current partially manual approach, with sufficient investment in computing and data storage infrastructure.

\section{Problems Encountered\label{sec:problems}}

The Spatial Aggregate (4km) uses a fixed grid support to group 5 or more observations in one grid cell, after which they are averaged to a single value. This was done to make the downstream data preparation and data matching steps more computationally efficient. Future work can conduct density based clustering like in \cite{yu2019high} (which is a more accurate representation of spatial grouping) but for high resolution predictors compared to the paper's 0.05 resolution. Spatial aggregation also reduced the observed SIF range, as there were only a few single footprint SIF samples in the $[0.75,1.5)$ range for the nine MGRS tiles in Germany, an expanded geographic scope is needed to obtain more very high SIF samples and therefore spatial aggregation can cover this range more reliably. As mentioned in \cite{doughty2022global}, both spatial and temporal aggregation can be done for OCO-2 SIF data to reduce SIF retrieval noise. But to make relatively high spatial resolution aggregates, temporal aggregation cannot work with OCO-2 data as over a period of 8-16 days, individual single date OCO-2 tracks can be separated by 10s of kilometres thereby increasing the spatial resolution of this aggregation. It would work for large scale studies of natural ecosystems or regional and state wide agriculture, but its coarse spatial resolution limits its usefulness for field level SIF reconstruction and analysis. Combining several OCO-2 dates could also mix different weather conditions and short term vegetation responses. Another issue was the temporal mismatch between SIF and Sentinel-2 optical data. OCO-2 SIF is attached to an individual acquisition date, VIIRS PAR is a daily mean, GLASS FAPAR is an eight-day composite, and Sentinel-2 WASP reflectance represents a monthly composite. The predictors were matched as closely as their temporal definitions allowed: PAR by exact date, FAPAR by the eight-day interval containing the SIF date and Sentinel-2 by its monthly product interval. Therefore the Sentinel-2 data represents more of an average monthly reflectance predictor data and thus could impede instantaneous SIF reconstruction. Native OCO-2 validation as seen in Figure \ref{fig:rq1_observed_predicted_comparison} can be noisy because the observed SIF may contain short term physiological responses that are not represented in the temporally aggregated Sentinel-2 predictors (monthly composites). 

Many of the previous SIF enhancement studies used MODIS predictors because they provided regular global composites at 250m, 500m resolution making data storage and processing much easier than handling 10m, 20m Sentinel-2 predictors. Moving from 500 m to 20 m increases the number of pixels covering the same area by a factor of 625. This therefore limited the spatial and temporal scope for acquiring Sentinel-2 data, eventhough SIF observation data for Germany ranged to about approximately 530k observations for years 2019-2024, the limited scope of Sentinel-2 data  reduced it to approximately 80k single footprint observations. Increasing the geographical scope could also add more samples from the low and high ends of the SIF distribution and could reduce the model's tendency to predict values near the centre of the training range. Another issue was that annual crop type rasters were themselves produced using Random Forest classification and spatial temporal filtering of Sentinel-1 and Sentinel-2 observations. They are therefore model outputs rather than error free original field records, \cite{gessner_crop_2025} reported a winter wheat F1-score of 0.82 and overall crop map accuracy of 75.5\%. Therefore classification errors from this data can propagate into the SIF model because they were used as crop composition predictors. The V11r OCO-2 SIF data used for this thesis does not include the correction for rotational Raman scattering (helps reduce SIF retrieval bias) described by Sanghavi et al. \cite{sanghavi2025impact} which will be available in future V12 data versions. Given the large inference times for generating 20m SIF maps, the yield experiment was conducted on only five MGRS tiles in Bavaria containing 48 NUTS 3 regions generating a dataset of 245 region-year records. Although the SIF maps were generated at 20 m resolution, the winter wheat pixels were averaged to one monthly value for each NUTS 3 region (limited by scope of yield data). Therefore much of the field scale spatial information was averaged out before yield prediction. Around roughly 290 NUTS 3 regions are available Germany wide (for harvest yield) potentially generating approximately 1,450 records, it would require extensive computational resources. Although these problems were not fully resolved within this thesis, they can be integrated in the current workflow with the addition of extra compute resources.


\section{Outlook\label{sec:outlook}}

In conclusion, future work should focus on obtaining denser and better corrected SIF observations, improved temporal alignment between SIF observations and predictor variables and evaluate the resulting enhanced SIF data at a spatial scale closer to their 20m output resolution. The European Space Agency's (ESA) FLEX mission \cite{coppo2017fluorescence, drusch2016fluorescence}, which is scheduled for launch in September 2026, will provide multi band SIF observations at approximately 300m resolution with more spatially and temporally complete samples. This 300m dataset can be the target support for future SIF enhancement downscaling methods or provide validation support for past SIF enhancement methods in Table \ref{tab:sif_enhancement_literature} including the thesis's method. Instead of using OCO-2 as the target support, TROPOMI SIF data can be used along with Sentinel-2 data for downscaling as it can include parent magnitude support during model training. As noted by \cite{sanghavi2025impact}, Build 12 of OCO-2/3 Level-2 products will provide a more stable bias corrected version of current OCO-2/3 products and can be used for the current thesis workflow when it is made available. With the addition of improved compute \& storage resources, cloud free Sentinel-2 Level 2A rasters closer to each SIF observation date can be used instead of monthly WASP Sentinel-2 composites as they will provide a much better temporal alignment between SIF and Sentinel-2 predictors. For the agricultural application, confidential parcel level yield observations (versus current NUTS 3 / district level) should be requested from the respective authorities, similar to the reference data used by Brandt et al. \cite{brandt2024ensemble} as it will greatly increase the size of the training data and give better understanding of the differences between crop pure and mixed SIF signals when done at field scales. Together, these additions could improve the accuracy, spatial detail and agricultural relevance of future enhanced SIF products.